\par\noindent\textbf{(1)} Maxwell静态边界满足切向E连续，磁场跳跃为 n̂×(H\_+−H\_−)=K\_s。二维Hall本构 K\_s=σ\_H ẑ×E\_t（符号随界面法线和Hall定义选取），故电荷激发的径向E在表面产生方位角环流，而非纵向Ohm电流。
\par\noindent\textbf{(2)} 弱Hall近似可先用无Hall静电界面点电荷场：若介电率上下为 ε₁,ε₂，界面径向电场 E\_r(r,0)=Q r/[2π(ε₁+ε₂)(r²+d²)\textasciicircum{}\{3/2\}]。从而 K\_φ(r)=σ\_H E\_r(r,0)，其方向依σ\_H符号；正反Hall响应翻转磁环流。
\par\noindent\textbf{(3)} 任意半径环带 dr 的电流 dI=K\_φ(r)dr，磁偶极矩 d m\_z=πr² dI。所以 m\_z(R\_c)=πσ\_H∫₀\textasciicircum{}\{R\_c\}r²E\_r(r,0)dr。在r≫d有 E\_r∼const/r²，因此被积量趋常数，整体磁矩随截止R\_c线性发散：无限平面电流尾巴不能归结为有限、局域磁偶极，必须保留完整空间分布或有限样品边界。
\par\noindent\textbf{(4)} 通过满足边界的磁标势形式，有时可在某一半空间把界面磁响应表示为位于镜像位置的形式磁电荷；它只是解边值问题的数学源，并不意味着观测到磁单极粒子。σ\_H→0 时所有Hall表面电流及相关磁响应消失；真实材料还需考虑非局域电导、有限频率与样品尺寸。
